\centerline{\bf \Huge Памятка по линейным движениям}

\textbf{Определение.}
Точка (прямая, вектор, что угодно) $A$ \textit{движется линейно}, если существует такой вектор $\bar{v}$, что $A(t)=A(0)+t \cdot \bar{v}$, где $A(0)$ это начальное положение, а $\bar{v}$ это так называемый вектор скорости вдоль которого движется наш объект. Нетрудно понять, что для точек и прямых линейное движение это просто параллельный перенос.

\href{https://www.geogebra.org/calculator/aqc37grt}{Примеры линейного движения (можно нажать)}


\problem{1}
\textbf{Основные свойства линейного движения. Их надо запомнить!}

\begin{itemize}
    \item[1)] Прямая постоянного направления, проведённая через линейно движущуюся точку, движется линейно.
    \item[2)] Середина отрезка (\textbf{да и в целом любая точка, делящая в фиксированном отношении!}), соединяющего две линейно движущиеся точки, движется линейно.
    \item[3)] Точка пересечения линейно движущихся прямых движется линейно.
\end{itemize}

\problem{2}
\textbf{Очень важное и полезное утверждение! Запомнить!}

Если зафиксировать вершины $A$ и $B$ треугольника $ABC$,а точку $C$ линейно двигать по прямой $AC$, то тогда центр описанной окружности, точка пересечения высот и точка пересечения медиан треугольника $A B C$ движутся линейно.

\textbf{Олимпиадный сленг: }Далее про два разных значения параметра $t$ при линейном движении объекта $A(t)$ будем говорить <<два разных положения>>

\problem{4}
\textbf{Проверка того, что точка лежит на какой-то прямой.} Если линейно движущаяся точка лежит на данной прямой в двух положениях, то она всегда на ней лежит.

\problem{5}\textbf{Проверка того, что три прямые пересекаются в одной точке.}
Если три линейно движущиеся прямые пересекаются в одной точке при двух положениях, то они всегда пересекаются в одной точке.

\problem{6}\textbf{Проверка перпендикулярности/параллельности двух прямых.} Если две линейно движущиеся прямые перпендикулярны/параллельны при трёх положениях, то они всегда перпендикулярны/параллельны.

\problem{7}\textbf{Проверка того, что три точки лежат на одной прямой.} Если три линейно движущиеся точки коллинеарны при трёх положениях, то они всегда коллинеарны.


\newpage

\centerline{\bf \Huge }

\textbf{В каких случаях полезно применять движение?}

--- утверждения, связанные с наличием равных отрезков, которые находятся далеко друг от друга (и не образуют простых конфигураций типа средней линии). В таком случае полезно подвигать концы этих отрезков или точки на них (например, середины);

--- утверждения, связанные с пересечением трех прямых в одной точке и с попаданием трех точек на одну прямую (правила 5 и 7), а также параллельностью или перпендикулярностью двух прямых (правило 6).


\textbf{Показания к применению}

--- не нужно двигать «стартовые» объекты, такие, как вершины исходного треугольника, потому что тогда вместе с ними поедет и вся конфигурация, и следить за ней будет тяжело;

--- не двигать жестко закрепленные объекты, особенно связанные с окружностями (поскольку окружность - объект нелинейный);

--- часто бывает полезно двигать точки в такие положения, которые вырождают конфигурацию; обычно этого удается добиться, склеив движущуюся точку с какой-то неподвижной;

--- иногда полезно заставить точки двигаться так, чтобы их скорости были пропорциональны (а не равны). Часто нет необходимости даже задумываться над точным значением коэффициента пропорциональности, достаточно лишь правильно описать само движение.

Особенно следует отметить, что иногда жесткость некоторых объектов в условии задачи оказывается мнимой! Т.е. утверждение задачи справедливо в гораздо большей общности, нежели той, которая предлагается вам. Сделано это, как правило, с целью запутать решателя и скрыть истинные идеи решения.


\newpage

\centerline{\bf \Huge }

\textbf{Несколько примеров решения задач с помощью линейного движения.}

Перед прочтением решения задачи обязательно сделайте у себя рисунок!

\problem{1}
Вписанная в треугольник $A B C$ окружность касается его сторон $A B$ и $A C$ в точках $C_1$ и $B_1$ соответственно. На отрезках $B C_1$ и $A B_1$ отмечены точки $P$ и $Q$, такие, что $P C_1=Q B_1$. Докажите, что середина отрезка $P Q$ лежит на прямой $B_1 C_1$.

\textbf{Решение.} Пусть точки $P$ и $Q$ линейно движутся из точек $C_1$ и $B_1$ с одинаковыми скоростями. Согласно правилу 7 (на самом деле можно и с помощью правила 4), нам необходимо найти три момента времени, когда середина отрезка $P Q$ лежит на отрезке $B_1 C_1$. Первый момент времени --- стартовый, т.е. $P=C_1$ и $Q=B_1$, второй --- когда $Q=A$ (тогда $Q C_1=Q B_1=P C_1$ и точка $C_1-$ середина $P Q$ ), а третий --- когда $P=A$ (тогда $P B_1=P C_1=Q B_1$ и точка $B_1-$ середина $P Q$ ). Обратите внимание, что в процессе доказательства точки $P$ и $Q$ выехали за пределы отведенных им отрезков $B C_1$ и $A B_1$ ! Это и есть мнимая жесткость условия: на самом деле его можно усилить, причем это усиление и приводит к решению!

\problem{2}
$\bigl[\textbf{ММО 2009, 10.5}\bigr]$
Стороны $B C$ и $A C$ треугольника $A B C$ касаются соответствующих вневписанных окружностей в точках $A_1$ и $B_1$. Пусть $A_2, B_2$ - ортоцентры треугольников $C A A_1$ и $C B B_1$. Докажите, что прямая $A_2 B_2$ перпендикулярна биссектрисе угла $C$.

\textbf{Решение.}
Ключевым моментом в этой задаче является соображение о том, что вневписанные окружности, грубо говоря, не при чем. Более точно, утверждение задачи остается справедливым для любых точек $A_1, B_1$ на сторонах $C A$ и $C B$, для которых $A B_1=B A_1$ (в случае, когда $A_1$ и $B_1$ - это точки касания, это равенство следует из свойств отрезков касательных). Поэтому в данной задаче нужно линейно двигать точки $A_1$ и $B_1$ из вершин $B$ и $A$ с одинаковыми по модулю скоростями. Обратите внимание, опять мнимая жесткость! Очевидно, что при таком движении точек $A_1$ и $B_1$ ортоцентры $A_2$ и $B_2$ будут двигаться линейно, причем их стартовая позиция - это ортоцентр $H$ треугольника $A B C$, а прямые, по которым они двигаются - это в точности высоты $A H$ и $B H$. Однако нам потребуется более сильное утверждение: ортоцентры $A_2$ и $B_2$ двигаются с одинаковыми по модулю скоростями. Это ожидаемо, поскольку как иначе использовать равенство скоростей точек $A_1$ и $B_1$ ?

\newpage

\hspace{3mm}

Доказать это легче всего простым счетом. Обозначим через $x$ длину отрезка $A_1 B$. Тогда длина отрезка $H A_2$ равна $x \operatorname{ctg} \angle B A C$ и не зависит только от длин сторон треугольника $C A A_1$. Поэтому длина отрезка $H B_2$ будет точно такой же, т.к. $A_1 B=B_1 A$.

Теперь самое время понять, при чем здесь биссектриса угла $\angle A C B$ и почему она должна быть перпендикулярна прямой $A_2 B_2$. Возможно, сходу это не очень понятно, но зато можно указать другую биссектрису, которая перпендикулярна прямой $A_2 B_2$. Это биссектриса треугольника $A_2 B_2 H$, который, как мы доказали, является равнобедренным. Поэтому осталось доказать, что биссектриса угла $\angle A H B$ параллельна биссектрисе угла $\angle A C B$. Мы полностью убрали из рассмотрения непонятные ортоцентры $A_2$ и $B_2$!

Дальше уже дело техники. Например, можно рассуждать так. Пусть $P$ и $Q$ - основания высот $A H$ и $B H$ соответственно. Рассмотрим четырехугольник $C P H Q$. Пусть биссектриса угла $\angle P C Q$ пересекает прямую $H P$ в точке $K$. Тогда

$$
\angle C K P=90^{\circ}-\angle P C Q / 2=90^{\circ}-\left(180^{\circ}-\angle P H Q\right) / 2=\angle P H Q / 2
$$

откуда и следует искомая параллельность. Задача решена!

\problem{3}
$\bigl[\textbf{IMO 2018}\bigr]$
На сторонах $A B$ и $A C$ треугольника $A B C$ выбраны точки $D$ и $E$, такие, что $A D=A E$. К отрезкам $B D$ и $C E$ проведены серединные перпендикуляры, пересекающие описанную окружность треугольника $A B C$ в точках $F$ и $G$. Докажите, что прямые $D E$ и $F G$ параллельны или совпадают.

\textbf{Решение.} Да-да задачи с межнара тоже вполне себе двигаются! Как мы сейчас увидим, рассуждения, необходимые в этой задаче, почти дословно повторяют решение задачи 2 с заменой ортоцентра на центр описанной окружности.

Ясно, что необходимо линейно двигать точки $D$ и $E$ из вершины $A$ с одинаковыми скоростями. Конечно, прямая $D E$ будет двигаться линейно. Однако линейность движения прямой $F G$ совершенно неочевидна (хотя она имеет место быть). Поэтому постараемся использовать связь наших точек с биссектрисой угла $\angle B A C$. Ясно, что эта биссектриса перпендикулярна $D E$. Если мы докажем, что эта же биссектриса перпендикулярна $F G$, то задача будет решена.


\newpage

\hspace{3mm}

Заметим, что если точки $D$ и $E$ линейно двигаются по прямым $A B$ и $A C$ с одинаковыми скоростями, то серединные перпендикуляры $p(t)$ и $q(t)$ к отрезкам $D B$ и $E C$ тоже будут двигаться линейно с одинаковыми скоростями, равными половинам скоростей точек $D$ и $E$. Рассмотрим точки $P(t)=p(t) \cap q(0)$ и $Q(t)=q(t) \cap p(0)$. Ясно, что точки $P(t)$ и $Q(t)$ двигаются линейно. Более того, рассуждая как в задаче $\mathbf{9}$, легко показать, что скорости точек $P(t)$ и $Q(t)$ равны. Это означает, что точка $O(t)$ пересечения $p(t)$ и $q(t)$ движется линейно по биссектрисе $\ell$ угла, образованного серединными перпендикулярами $p(0)$ и $q(0)$ к сторонам треугольника $A B C$, с вершиной в центре $O$ описанной вокруг него окружности. Ясно, что прямая $F G$ перпендикулярна $\ell$ в силу осевой симметрии относительно $\ell$.

Доказательство того, что прямая $\ell$ параллельна биссектрисе угла $\angle B A C$, абсолютно аналогично рассуждению из задачи 2. Это завершает решение.